\documentstyle[12pt,leqno,bezier]{cicfinal}

\newcommand{\mar}[1]{\marginpar[\raggedright\footnotesize\sf #1]{\raggedleft\footnotesize\sf #1}}
\newcommand{\asideleft}[1]{\medskip\noindent%
	\rule{321pt}{0pt}\makebox[0pt][r]{%
	\begin{minipage}[t]{400pt}\footnotesize\sf{#1}%
	\end{minipage}}\bigskip}
\newcommand{\asideright}[1]{\medskip\noindent%
	\rule{40pt}{0pt}\makebox[0pt][l]{%
	\begin{minipage}[t]{400pt}\footnotesize\sf{#1}%
	\end{minipage}}\bigskip}
\newcommand{\aside}[1]{\ifodd \value{page}%
	{\asideright{#1}}\else{\asideleft{#1}}\fi}
\newcounter{exercisenumber}[subsection]
\newcounter{exercisepart}[exercisenumber]
\newcommand{\num}{\stepcounter{exercisenumber}\par\bigskip%
	\noindent\arabic{exercisenumber}.\quad}
\newcommand{\numa}{\stepcounter{exercisenumber}\par\bigskip%
	\noindent\arabic{exercisenumber}.\quad%
	\stepcounter{exercisepart}\alph{exercisepart})\enskip}
\newcommand{\sub}{\stepcounter{exercisepart}\par\smallskip%
	\noindent\alph{exercisepart})\enskip\thinspace}
\newcommand{\ans}[1]{\par\smallskip\noindent[Answer:\enskip%
	{#1}]}
\newcommand{\dint}{\displaystyle \int}

\makeindex

\begin{document}
\setcounter{page}{503}
\setcounter{chapter}{8}



\chapter{Functions of Several Variables}

	Functions that depend on several input variables first appeared in the $S$-$I$-$R$ model at the beginning of the course.  Usually, the number of variables has not been an issue for us.  For instance, when we introduced the derivative in chapter 3, we used partial derivatives to treat functions of several variables in a parallel fashion.  
	\mar{The problem of visualizing a function of several variables}	
However, when there are questions of visualization and geometric understanding, the number of variables {\em does\/} matter.  Every variable adds a dimension to the problem---one way or another.  For example, if a function has two input variables instead of one, we will see that its graph is a surface rather than a curve. 
		
	This chapter deals with the geometry of functions of two or more variables.  We start with graphs and level sets.  These are the basic tools for visualization.  Then we turn to microscopic views, and see what form the microscope equation takes.  Finally, we consider optimization problems using both direct visual methods and dynamical systems.
	
	
	
	
\section{Graphs and Level Sets}

\special{psfile=../../ps/911/weather1.ps}

\begin{minipage}[t]{3.5in}
\quad \ The graph\index{graph} at the right comes from a model that describes how the average daily temperature\index{temperature function} at one place varies over the course of a year.  It shows the temperature $A$ in $^\circ$F, and the time $t$ in months from January. \hfill As we would expect, the
\end{minipage}  

\noindent
temperature is a periodic function (which we can write as $A(t)$), and its period is 12 months.  Furthermore, the lowest temperature occurs in February (when $t \approx$ 1 or 13) and the highest in July (when $t \approx$ 7 or 19).  This is about what we would expect.

	However, all these temperature fluctuations disappear 
	\mar{Underground temperatures fluctuate less}	
a few feet underground.  Below a depth of 6 or 8 feet, the temperature of the soil remains about 55$^\circ$F year-round!  Between ground level and that depth, the temperature still fluctuates, but the range from low to high decreases with the depth.  Here is what happens at some specific 
\linebreak
 

\special{psfile=../../ps/911/weather2.ps}
\vspace{166pt}
\label{timegraph}

\noindent
depths.  Notice how the time at which the temperature peaks gets later and later as we go farther and farther underground.  For example, at $d = 2$ feet the highest temperature occurs in September ($t \approx 9$), 
 	\mar{The phase shifts with the depth}
not July.  It literally takes time for the heat to sink in.  In chapter 7 we called this a phase shift.\index{phase shift}  The lowest temperature shifts in just the same way.  At a depth of 2 feet, it is colder in March than in January.   

	Thus $A$ is really a function of {\em two\/} variables,\index{function, of two variables} the depth $d$ as well as the time $t$.  To reflect this addition, let's change our notation for the function to $A(t, d)$.  In the figure above, $d$ plays the role of a parameter: it has a fixed value for each graph.  We can reverse these roles and make $t$ the parameter.  This is done in the figure on the top of the next page.  It shows us how the temperature 
	\mar{Graphing temperature as a function of depth}	
varies with the depth at fixed times of year. Notice that, in April and October, the extreme temperature is not found on the surface.  In October, for example, the soil is warmest at a depth of about 9 inches.
	
	The lower figure on the page is a single graph that combines all the information in these two sets of graphs.  Each point on the bottom 
\linebreak

\mbox{}

\special{psfile=../../ps/911/weather3.ps}
\vspace{150pt}
\label{depthgraph}
	
\noindent
of the box corresponds to a particular depth and a particular time.  The height of the surface above that point tells us the temperature at that depth and time.\index{graph, reading a}\index{surface graph}  
	\mar{Reading a surface graph}
For example, suppose you want to find the temperature 4 feet below the surface at the beginning of July.  Working from the bottom front corner of the box, move 4 feet to the right and then 6 months toward the back.  This is the point $(d, s) = (4, 6)$.  The height of the graph above this point is the temperature $A$ that we want.
    	
\special{psfile=../../ps/911/weather5.ps}
	
\newpage

	There is a definite connection between this surface and the two collections of curves.  
	\mar{Grid lines are slices of the surface that show how the function depends on each variable separately}	
Imagine that the box containing the surface graph is a loaf of bread.  If you slice the loaf parallel to the left or right side, this slice is taken at a fixed depth.  The cut face of the slice will look like one of the graphs on page \pageref{timegraph}.  These show how the temperature depends on the time at fixed depths.  If you slice the loaf the other way---parallel to the front or back face---then the time is fixed.  The cut face will look like one of the graphs on page \pageref{depthgraph}.  They show how the temperature depends on the depth at fixed times.  The grid lines\index{grid lines} on the surface are precisely these ``slice''\index{slice}\index{graph, slice of a} marks. 


\special{psfile=../../ps/911/weather6.ps}
\hfill
\begin{minipage}[t]{2.7in}
\quad \ Here is the same graph seen from a different viewpoint.  Now time is measured from the left, while depth is measured from the back.  The temperature is still the height, though.  One advantage of this view is that it shows more clearly how the peak temperature is phase-shifted with the depth.

\quad \ We now have two ways to visualize how the average daily temperature depends on the time of year and the depth below ground.  One is the surface graph itself, and the other is a collection of curves that are slices of the surface.  The surface gives us 
\linebreak
\end{minipage}
\vspace{-9pt}

\noindent
an overall view, but it is not so easy to read the surface graph to determine the temperature at a specific time and depth.   
	\mar{Comparing the surface to its slices}
Check this yourself: what is the temperature 2 feet underground at the beginning of April?  The slices are much more helpful here.  You should be able to read from either collection of slices that $A \approx 44^\circ$F.  

\subsection{Examples of Graphs}

\special{psfile=../../ps/911/minmax1.ps}

	The purpose of this section is to get some experience constructing and interpreting surface graphs.  To work in a context, look first at the functions $y = x^2$ and $y = -x^2$.  They provide us with standard examples of a minimum and a maximum when there is just one input variable.  Let's consider now the corresponding examples for two input variables.  Besides an ordinary maximum and an ordinary minimum, we will find a {\em third\/} type---called a minimax---that is completely new.  It arises because a function can have a minimum with respect to one of its input variables and a maximum with respect to the other.  
	
\subsubsection{A minimum: $z = x^2 + y^2$}
\index{minimum}

\special{psfile=../../ps/911/min2.ps}

\begin{minipage}[t]{3.7in}
\quad \	At the origin $(x, y) =(0, 0)$, $z = 0$.  At any other point, either $x$ or $y$ is non-zero. Its square is positive, so $z > 0$.  Consequently, $z$ has a minimum at the origin.  The graph of this function is a parabolic {\bf bowl} whose lowest point sits on the origin.  As always, the grid lines are slices, made by fixing the value of $x$ or $y$.  For example, if $y = c$, then the slice is $z = x^2 + c^2$.  This is an ordinary parabolic curve. 
\end{minipage}
		
		
\subsubsection{A maximum: $z = -x^2 - y^2$}
\index{maximum}

\special{psfile=../../ps/911/max2.ps}

\begin{minipage}[t]{3.7in}
	For any $x$ and $y$, the value of $z$ in this example is the opposite of its value in the previous one.  Thus, $z$ is everywhere negative, except at the origin, where its value is 0.  Thus $z$ has a maximum at the origin.  Its graph is an upside-down bowl, or {\bf peak}, whose highest point reaches up and touches the origin.  Grid lines are the curves $z = -x^2 - c^2$ and $z = -c^2 - y^2$.  These are parabolic curves that open downward.
\end{minipage}

\subsubsection{A minimax: $z = x^2 - y^2$}
\index{minimax}

\special{psfile=../../ps/911/saddle1.ps}

\label{saddle}
\begin{minipage}[t]{3.7in}
\quad \ Suppose we fix $y$ at $y = 0$.  This slice has the equation $z = x^2$, so it is an ordinary parabola (that opens upward).  Thus, as far as the input $x$ is concerned, $z$ has a {\em minimum\/} at the origin.  Suppose, instead, that we fix $x$ at $x = 0$. Then we get a slice whose equation is $z = -y^2$.  It is also an ordinary parabola, but this one opens downward.  As far as $y$ is concerned, $z$ has a {\em maximum\/} at the origin.  It is 
\linebreak
\end{minipage}
\vspace{-9pt}

\noindent
clear from the  graph how upward-opening slices in the $x$-direction fit together with downward-opening slices in the $y$-direction.  Because of the shape of the surface, a minimax is commonly called a {\bf saddle}, or a {\bf saddle point}.\index{saddle}\index{saddle point}

\newpage

	Here are two slices\index{slice}\index{graph, slice of a} of $z = x^2 - y^2$ shown in more detail.  Points in the box have three coordinates: $(x, y, z)$.  If we set $y = 0$ we are selecting 
\linebreak

\special{psfile=../../ps/911/saddle2.ps}

\special{psfile=../../ps/911/saddle3.ps}

\special{psfile=../../ps/911/saddle4.ps}
\vspace{200pt}

\noindent
\begin{minipage}{3.6in}
the points of the form $(x, 0, z)$.  These make up the $x, z$-plane.  On this plane the equation $z = x^2 - y^2$ becomes simply $z = x^2$.  The graph of {\em this\/} equation is a curve in the $x, z$-plane---specifically, the parabola shown.  The situation is similar if $y$ is given some other fixed  
\linebreak 
\end{minipage}
\vspace{-9pt}

\noindent
value. For example, $y = -4$ specifies the points $(x, -4, z)$.   
	\mar{\vspace{-36pt} The points where $y = c$ form a vertical plane parallel to the $x, z$-plane}
These describe the plane that forms the front face of the box.  The equation $z = x^2 - y^2$ becomes $z = x^2 - 16$.  The curve tracing out the intersection of the saddle with the front of the box is precisely the graph of $z = x^2 - 16$.  

	If $x = 0$ we get the points $(0, y, z)$ that make up the $y, z$-plane.  On this plane the equation simplifies to $z = -y^2$, and its graph is the parabolic curve shown.  
	\mar{The points where $x = c$ form a vertical plane parallel to the $y, z$-plane}
Giving $x$ a different fixed value leads to similar results.  A good example is $x = -4$.  The points $(-4, y, z)$ lie on the plane that forms the left side of the box.  The equation becomes $z = 16 - y^2$ there, and this is the parabolic curve marking the intersection of the saddle with the left side of the box.

\aside{As you can see, it is valuable for you to be able to generate surface graphs yourself.  There are now a number of computer utilities which will do the job.  Some can even rotate the surface while you watch, or give you a stereo view.  However, even without one of these powerful utilities, you should try to generate the slicing curves that make up the grid lines of the surface.}

\subsubsection{A cubic: $z = x^3 - 4x - y^2$}
\label{cubic}

	Slices of this graph are downward-opening parabolas (when $x = c$) and are cubic curves that have the same shape (when $y = c$).  Notice that each cubic curve has a maximum and a minimum, and each parabola 
\linebreak    

\special{psfile=../../ps/911/cubic1.ps}
\special{psfile=../../ps/911/cubic2.ps}
\vspace{110pt}

\noindent
has a maximum.  
	\mar{The surface has a saddle}
The surface graph itself has a {\em peak\/} where the cubics have their maximum, but it has a {\em saddle\/}\index{saddle} where the cubics have a minimum.  Do you see why?  The saddle point is a minimax for $z = x^3 -4x -y^2$: $z$ has a minimum there {\em as a function of $x$ alone} but a maximum {\em as a function of $y$ alone}.

\special{psfile=../../ps/911/cubic3.ps}
\special{psfile=../../ps/911/cubic4.ps}
\special{psfile=../../ps/911/cubic5.ps}
\vspace{230pt}

	The small figures on the right show the same surface as the large figure; they just show it from different viewpoints.  
	\mar{See the graph from different viewpoints}	
As a practical matter, you should look at these surfaces the way you would look at sculpture: ``walk around them'' by generating diverse views.   
	
\subsubsection{Energy of the pendulum: $E = 1 - \cos{\theta} + {1 \over 2} v^2$}

\special{psfile=../../ps/911/penden1.ps}
\vspace{200pt}

\special{psfile=../../ps/911/penden2.ps}
\hfill
\begin{minipage}[t]{3.5in}
\quad \ This function first came up in chapter 7, where it was used to demonstrate that a dynamical system describing the motion of a frictionless pendulum had periodic solutions.\index{pendulum}\index{energy}  It was used again in chapter 8 to clarify the phase portrait of that dynamical system.  The function $E$ varies periodically with $\theta$, and you can see this in the graph.  The minimum\index{minimum}\index{energy minimum} at the origin is repeated at $(\theta, v) = (2\pi, 0)$, and so on.  The graph also has a saddle at the point $(\theta, v) = (-\pi, 0)$.  This too repeats with period $2\pi$ in the $\theta$ direction.

\quad \ The figure at the left is the same surface with part cut away by a slice of the form $v = c$.  These slices are sine curves: $E = 1 - \cos{\theta} + {1 \over 2} c^2$.  Slices of the form $\theta = c$ are upward-opening parabolas.  From this viewpoint, the saddle points show up clearly.  
\end{minipage}

\aside{One way to describe what happens to a real pendulum---that is, one governed by frictional forces as well as gravity---is to say that its energy ``runs down'' over time.  Now, at any moment the pendulum's energy is a point on this graph. As the energy runs down, that point must work its way down the graph.  Ultimately, it must reach the bottom of the graph---the minimum energy point at the origin $(\theta, v) = (0, 0)$.  This is the stable equilibrium point.  The pendulum hangs straight down ($\theta = 0$) and is motionless ($v = 0$).  The graph gives us an abstract---but still vivid and concrete---way of thinking of the dissipation of energy.\index{dissipation of energy}}    
	

\subsection{From Graphs to Levels}

\special{psfile=../../ps/912/density2.ps}
\special{psfile=../../ps/912/density1.ps}

\begin{minipage}[t]{3.5in}
	There is still another way to picture a function of two variables.  To see how it works we can start with an ordinary graph.  On the right is the graph of 
	$$
	z = f(x, y) = x^3 - 4x - y^2,
	$$
the cubic function we considered on page~\pageref{cubic}.  This graph looks different, though.  The difference is that points are shaded according to their height.  Points at the bottom are lightest, points at the top are darkest.  

\quad \ Notice that the flat $x, y$-plane is shaded exactly like the graph above it.  For instance, the dark spot centered at the point $(x, y) = (-1, 0)$ is directly under the peak on the graph.  The other 
\linebreak 
\end{minipage}
\vspace{-9pt}

\noindent
dark patch, near the right edge of the plane, is under the highest visible part of the surface.  Consequently, the shading on the $x, y$-plane gives us the same information as the graph.  In other words, {\em the intensity of shading at $(x, y)$ is proportional to the value of the function $f(x, y)$}.

	The figure in the $x, y$-plane is called a {\bf density plot}.\index{density plot}  
	\mar{Density plots}	
Think of the intensity of shading as the {\em density\/} of ink on the page.  Here are density plots of the standard minimum,\index{minimum} maximum,\index{maximum} and minimax.\index{minimax}  Compare
\label{density}
\linebreak
	
\special{psfile=../../ps/912/density3.ps}
\vspace{140pt}

\noindent
\mbox{
	\small
	\mbox{} \hspace{25pt}
	$z = x^2 + y^2$
	\hspace{70pt}
	$z = -x^2 - y^2$
	\hspace{80pt}
	$z = x^2 - y^2$
}

\vspace{6pt}

\noindent
these with the graphs on page~\pageref{saddle}.  The third density plot is the most interesting.  From the center of the $x, y$-plane, the shading increases to the right and left.  Therefore, $z$ has a minimum in the horizontal direction.  However, the shading {\em decreases\/} above and below the center.  Therefore, $z$ has a {\em maximum\/} in the vertical direction.  Thus, you really can see there is a minimax at the origin.

	Try your hand at reading the density plot on the left below.  
	\mar{A sample plot}	
You should see two maxima (directly above and below the origin), a minimum (at the origin itself), and two saddles (to the right and the left of the origin).  The function defining the plot is
	$$
	f(x, y) = (x^2 + (y-1)^2 - 3)(3 - x^2 - (y+1)^2).
	$$

\special{psfile=../../ps/912/density4.ps}
\special{psfile=../../ps/912/density5.ps}
\vspace{210pt}

\noindent
Can you visualize what the graph looks like?  This density plot should help you, and you can also construct slices by setting $x = c$ and $y = c$.  The slices $x = 0$ and $y = 0$ are especially useful.  With them you could determine the exact coordinates of the maxima and the saddles.

\aside{These density plots show a ``checkerboard'' pattern because {\em Mathematica\/}\index{Mathematica} (the computer program that produces them) shades each little square according to the value of the function at the center of the square.  This pattern is an artefact; it is not inherent to a density plot.} 

	In a density plot, the shading varies smoothly with the value of the function.  This is accurate, but it may be a bit difficult to read.  On the right you see a modified density plot.  There is still shading, but there are now just a few distinct shades.  
	\mar{From densities \\to contours}
This makes a sharp boundary between one shade and the next.  The boundary is called a {\bf contour},\index{contour} or a {\bf level}.\index{level}  The figure itself is called a {\bf contour plot}.\index{contour plot}  The two maxima on the vertical line $x = 0$ stand out more clearly on the contour plot.  Also, the contour lines around the two saddles help us see that the function has a minimum in the vertical direction and a maximum in the horizontal direction. 
 

\special{psfile=../../ps/913/contour1.ps}  

\noindent
\begin{minipage}[t]{3.4in}
\quad \ Once we have contour lines to separate one density level from the next, we can even dispense with the shading.  The figure on the right is just the contour plot from the opposite page, minus the shading.  The contour lines, or level curves, now stand out clearly.  On each contour, the value of the function is constant.  This is also called a {\bf contour plot}.  

\quad \ There is some loss of information here, however.  For example, we can't tell where the value of the function is large and where it is small.  Nevertheless, the nested ovals on the vertical line $x = 0$ {\em do\/} tell us that there is either a maximum or a minimum at the center of each nest. 
\end{minipage}
\vspace{5pt}

	For reference purposes, here are the contour plots for the standard minimum,\index{minimum} maximum,\index{maximum} and saddle.\index{saddle}  In the first two cases, the contours are concentric ovals.  
	\mar{Contours of the standard functions}	
These look the same, so only one is illustrated.  The other two pictures show a saddle.  In general, the contours around a saddle are a family of hyperbolas.  However, it is possible for one of the contour lines to pass exactly through the minimax point.  That contour is a pair of crossed lines, as shown in the version on the right.  You should compare these contour plots with the density plots of the same functions on page~\pageref{density}, and with their graphs on page~\pageref{saddle}.
	
\special{psfile=../../ps/913/contour2.ps}  
\special{psfile=../../ps/913/contour3.ps}  
\special{psfile=../../ps/913/contour4.ps}  

 \vspace{160pt}

{\small
\begin{picture}(360,10)
	\put(52,12){\makebox(0,0){two functions but one plot}}
	\put(52,-2){\makebox(0,0){$z = x^2 + y^2$}}
	\put(52,-16){\makebox(0,0){$z = -x^2 - y^2$}}
	\put(283,8){\makebox(0,0){two plots of a single function}}
	\put(283,-8){\makebox(0,0){$z = x^2 - y^2$}}
\end{picture}
}

\newpage

	There is a direct connection between the contour plot 
	\mar{Contours are horizontal slices of a graph}
of a function and its graph.\index{graph, slices of a}\index{horizontal slice of a graph}  Contours are horizontal slices of the graph, just as grid lines are vertical slices.  Below, we use the standard functions $z = x^2 - y^2$ and $z = x^2 + y^2$ to illustrate the connection.  Notice that every contour down in the $x, y$-plane lies exactly below, and has the same shape as, a horizontal slice of the graph.  This picture explains why contours are called {\em level\/} curves.\index{level curve}  
	
\special{psfile=../../ps/913/consurf1.ps}
\special{psfile=../../ps/913/consurf2.ps}
\vspace{190pt}


	To get some more experience with contour plots, we return to 
	\mar{Energy of the pendulum, \\again}	
the energy\index{pendulum}\index{energy} function of the pendulum:
	$$
	E(\theta, v) = 1 - \cos{\theta} + {\textstyle {1 \over 2}} v^2.
	$$
From the contour plot alone you should be able to see that $E$ has either a minimum or a maximum at $(\theta, v) = (0, 0)$, and another at $(2\pi, 0)$.  The contours also provide evidence that there is a saddle (minimax) near $(\theta, v) = (-\pi, 0)$ and $(\pi, 0)$.  It is also apparent that $E$ is a {\em periodic\/} function of $\theta$.  

\special{psfile=../../ps/913/contour5.ps}

\newpage

	What you should find most striking about this plot, however, is the way it resembles the phase portrait of the pendulum (chapter 8, pages 464--467).  
% qqq
Every level curve here looks like a trajectory\index{trajectory} of the dynamical system.\index{dynamical system}  
	\mar{Energy contours are trajectories of the dynamics}	
This is no accident.  We know from chapter 8 that the energy is a first integral for the dynamics.\index{first integral}  In other words, energy is constant along each trajectory---this is the law of conservation of energy.  But each level curve shows where the energy function has some fixed value.  Therefore, each trajectory must lie on a single energy level. 

\special{psfile=../../ps/913/contour6.ps}

\noindent
\begin{minipage}[t]{3.3in}
\quad \ We can carry the connection between contours and trajectories even further.  Closed trajectories correspond to {\em oscillations\/}\index{oscillation} of the pendulum.  But the closed trajectories are the closed contours, and these are the ones that surround the minimum.  In particular, they are {\em low\/} energy levels.  By contrast, at higher energies ($E > 2$, in fact), the pendulum will just continue to spin in what ever direction it was moving initially.  Thus, each high energy level is occupied by {\em two\/} trajectories---one for clockwise spinning and one for counter-clockwise.
\end{minipage}


\subsection{Technical Summary}

	The examples we have seen so far were meant to introduce some of the common ways of visualizing a function $z = f(x, y)$.  To use them most effectively, though, you need to know more precisely how each is defined.  We review here the definition of a graph, a density plot, a contour plot, and a terraced density plot.  
	
\special{psfile=../../ps/914/sumgraph.ps}

\medskip
\noindent
\begin{minipage}[t]{3.3in}
{\bf Graph}.\index{graph, surface}  The graph of $z = f(x, y)$ lies in the 3-dimensional space with coordinates $(x, y, z)$.  To construct it, take any input $(x, y)$.  Identify this with the point $(x, y, 0)$ in the $x, y$-plane (which is defined by the condition $z = 0$).  The corresponding point on the graph lies at the height $z = f(x, y)$ above the $x, y$-plane.  This point has coordinates $(x, y, f(x, y))$.   {\bf The graph is the set of all points of the form $(x, y, f(x, y))$.}  This is a 2-dimensional surface.     

\end{minipage} 

\newpage

\special{psfile=../../ps/914/sumdense.ps}
\special{psfile=../../ps/914/sumcon2.ps}
\special{psfile=../../ps/914/sumcon4.ps}
\hfill
\begin{minipage}[t]{3.7in}
{\bf Density plot}.\index{density plot}  The density plot of $z = f(x, y)$ lies in the 2-dimensional $x, y$-plane.  Choose any rectangle where the function is defined, and let $m$ and $M$ be the minimum and maximum values, respectively, of $f(x, y)$ on the rectangle.  Define
	$$
	\rho(x, y) = {f(x, y) - m \over M - m}.
	$$
Then $\rho$ satisfies $0 \leq \rho(x, y) \leq 1$ on the rectangle; it is called a {\bf density function} ($\rho$ is the Greek letter {\em rho\/}).  {\bf In the density plot, the density of ink---or darkness---at $(x, y)$ is $\rho(x, y)$.}
\medskip

{\bf Contour plot}.\index{contour plot}\index{contour}  A {\bf contour} of $z = f(x, y)$ is the set of points in the $x, y$-plane where $f$ has some fixed value: 
	$$
	f(x, y) = c.
	$$
That fixed value $c$ is called the {\bf level}\index{level} of the contour.   (The two solid ovals in the figure at the left make up a single contour.)  A contour is also called a level curve.  {\bf A contour plot of $f$ is a collection of curves $f(x, y) = c_j$ in the $x, y$-plane.}  In the plot it is customary to use constants $c_1, c_2, \ldots$ that are equally spaced; that is, the interval between one $c_j$ and the next always has the same value $\Delta c$.
\medskip

{\bf Terraced density plot}.\index{density plot, terraced}\index{terraced density plot}  This is a contour plot in which the region between two adjacent contours is shaded with ink of a single density.  If the contours are at levels $c_1$ and $c_2$, then the density that is typically chosen is the one for the level half-way between these two---that is, for their average $(c_1 + c_2)/2$.    Each region is called a {\bf terrace}.\index{terrace}  Often, a terraced density plot is drawn in color, using different colors for each terrace.  Television weather programs use terraced density plots to describe the temperature forecast for a large region.
\end{minipage}

	\aside{We find density plots everywhere.  A photograph is a density plot of the light that fell on the film when it was exposed.  A newspaper `half-tone' illustration is also a density plot of an image.}

\subsubsection{The pros and cons}

	Each of these modes of visualization has advantages and disadvantages.  All are reasonably good at indicating the extremes (the maxima and minima) of a function.  A contour plot needs some additional information---for example, a label on each contour to indicate its level---to distinguish between maxima and minima.  However, if you want to know the numerical value of $f(x, y)$ at a particular point $(x, y)$, a contour plot with labels offers more precision than a density plot.  It's usually better than a graph, too.  

	Overall, a graph has the biggest visual impact, 
	\mar{Plots are visually economical in comparison to graphs}	
but there is a cost.  It takes three dimensions to represent the graph of a function of two variables, but only two to represent a plot.  The cost is that extra dimension.  It means that we cannot draw the graph of a function of three variables.  That would take four mutually perpendicular axes---an impossibility in our three-dimensional space.  However, we can produce a contour plot.  

\subsection{Contours of a Function of Three Variables}

	We pause here for a brief glimpse of a large subject.  By analogy with the definition for a function of two variables,
	\mar{Contours and levels}	
 we say that a {\bf contour}\index{contour} of the function $f(x, y, z)$\index{function, of three variables} is the set of points $(x, y, z)$ that satisfy the equation      
	$$
	f(x, y, z) = c,
	$$
for some fixed number $c$.  We call $c$ the {\bf level}\index{level} of the contour. 

	Let's find the contours of $w = x^2 + y^2 + z^2$.  
	\mar{The standard minimum}	
This is completely analogous to the function $x^2 + y^2$ with two inputs.  (What do the contours of $x^2 + y^2$ look like?)  In particular, $w$ has a minimum when $(x, y, z) = (0, 0, 0)$.\index{minimum}   As the following diagram shows, $w = x^2 + y^2 + z^2$ is the square of the distance from the origin to $(x, y, z)$.  (We use the Pythagorean theorem twice: once for $p^2$ and once for $q^2$.)  Consequently, 
\linebreak

\special{psfile=../../ps/915/xyz1.ps}
\vspace{15pt}
	   
	\hspace{180pt}
	\begin{tabular}{rcl}
	$p^2$ &=& $x^2 + y^2$ \\ [3pt]
	$q^2$ &=& $p^2 + z^2$ \\
	      &=& $x^2 + y^2 + z^2$ \\
	      &=& $w$
	\end{tabular}
	
\newpage

\noindent
all points $(x, y, z)$ where $w$ has a fixed value lie a fixed distance from the origin.  Specifically, $w = x^2 + y^2 + z^2 = c$ is the following set:
	\begin{itemize}
\item $c > 0$: the sphere of radius $\sqrt{c}$ entered at the origin;
\item $c = 0$: the origin itself;
\item $c < 0$: the empty set.
\end{itemize}
The contour plot of $w = x^2 + y^2 + z^2$ is thus 
	\mar{The contours \\are spheres}
a nest of concentric spheres, as shown in the illustration below.\index{level surface}   The value of $w$ is constant on each sphere.  (The tops of the spheres have been cut away so you can see how the spheres nest; the whole thing resembles an onion.) 

\special{psfile=../../ps/915/shells1.ps}

\newpage

	Below is the contour plot of another standard function 
	\mar{A standard minimax}	
with three input variables:
	$$
	w = f(x, y, z) = x^2 + y^2 - z^2.
	$$
A quarter of each surface has been cut away so you can see how the surfaces nest together.  Note that $w = 0$ is a cone, and every surface with $w < 0$ consists of two disconnected (but congruent) pieces---an upper half and a lower half.\index{minimax}

	You should compare this function to the standard minimax $x^2 - y^2$ in two variables.  The three-variable function $w$ has a minimum with respect to {\em both\/} of the variables $x$ and $y$, while it has a maximum with respect to $z$.  (Do you see why?  The arguments are exactly the same as they were for two input variables on page~\pageref{saddle}.)  
	\mar{The contours \\are hyperbolic shapes}	
Furthermore, the contours of $x^2 - y^2$ are a family of hyperbolas, and the contours of $x^2 + y^2 - z^2$ are surfaces obtained by rotating these hyperbolas about a common axis.

\special{psfile=../../ps/915/shells2.ps}
\special{psfile=../../ps/915/shells3.ps}

\newpage

	It is a general fact---and our two examples 
	\mar{When there are three input variables, the contours are surfaces}	
provide good evidence for it---that a single contour of a function of three variables is a {\em surface}.  Thus a contour is a curve or a surface, depending on the number of input variables.  We often use the term {\bf level set}\index{level set} (rather than a level {\em curve\/} or a level {\em surface\/})\index{level surface}\index{surface, level} as a generic name for a contour.  

\subsection{Exercises}

	In many of these exercises it will be essential to have a computer program to make graphs, terraced density plots, and contour plots of functions of two variables.
	
\numa  Use a computer to obtain a graph of the function $z = \sin{x} \sin{y}$ on the domain $0 \leq x \leq 2 \pi$, $0 \leq y \leq 2 \pi$.  How many maximum points do you see?  How many minimum points?  How many saddles?

\sub  Determine, as well as you can from the graph, the location of the maximum, minimum, and saddle points.

\num  Continuation.  Make the domain $-2 \pi \leq x \leq 2 \pi$, $-2 \pi \leq y \leq 2 \pi$ and answer the same questions you did in the previous exercise.  (Does the graph look like an egg carton?)

\num  Obtain a terraced density plot (or a contour plot) of $z = \sin{x} \sin{y}$ on the domain $-2 \pi \leq x \leq 2 \pi$, $-2 \pi \leq y \leq 2 \pi$.  Locate the maximum, minimum, and saddle points of the function.  Do these results agree with those from the previous exercise?

\num  Obtain the graph of $z = \sin{x} \cos{y}$ on the domain $0 \leq x \leq 2 \pi$, $0 \leq y \leq 2 \pi$.  How does this graph differ from the one in exercise~1?  In what ways is it similar?

\num  Obtain the graph of $z = 2x + 4 x^2 - x^4 - y^2$ when $-2 \leq x \leq 2$, $-4 \leq y \leq 4$.  Locate all the minimum, maximum, and saddle points in this domain.  [Note: the minimum is on the boundary!]

\num  Continuation.  Obtain a terraced density plot (or contour plot) for the function in the previous exercise, using the same domain.  Use the plot to locate all the minimum, maximum, and saddle points.  Compare your results with those of the previous exercise.

\numa  Obtain the graph of $z = 2x - y$ on the domain $-2 \leq x \leq 2$, $-2 \leq y \leq 2$.  What is the shape of the graph?

\sub Graph the same function of the domain $2 \leq x \leq 6$, $0 \leq y \leq 4$.  What is the shape of the graph?  How does this graph compare to the one in part (a)?

\numa Continuation.  Sketch three different slices of the graph of $z = 2x - y$ in the $y$-direction.  What do the slices have in common?  How are they different?

\sub  Answer the same questions for slices in the $x$-direction.

\numa  Obtain the graph of $z = .3 x + .8 y + 2.3$; choose the domain yourself.  
Where does the graph intercept the $z$-axis?

\sub  Describe the vertical slices of this graph in the $y$-direction and in the $x$-direction.

\num  Describe the vertical slices of the graph of $z = px + qy + r$ in the $y$-direction and in the $x$-direction.

\numa  Compare the contours of the function $z = x^2 + 2 y^2$ to those of $z = x^2 + y^2$.  

\sub  What is the shape of the graph of $z = x^2 + 2 y^2$?  Decide this first using only the information you have about the contours.  Then use a computer to obtain the graph.

\numa  Compare the contours of the function $z = x^2 - 2 y^2$ to those of $z = x^2 - y^2$.  

\sub  What is the shape of the graph of $z = x^2 - 2 y^2$?  Decide this first using only the information you have about the contours.  Then use a computer to obtain the graph.

\numa  Obtain a contour plot of the function $z = x^2 + xy + y^2$.  

\sub  What is the shape of the graph of $z = x^2 + xy + y^2$?  Decide this first using only the information you have about the contours.  Then use a computer to obtain the graph.

\numa  Obtain a contour plot of the function $z = x^2 + 3xy + y^2$.  

\sub  What is the shape of the graph of $z = x^2 + 3xy + y^2$?  Decide this first using only the information you have about the contours.  Then use a computer to obtain the graph.

\numa  Obtain a contour plot of the function $z = x^2 + 2xy + y^2$.  

\sub  What is the shape of the graph of $z = x^2 + 2xy + y^2$?  Decide this first using only the information you have about the contours.  Then use a computer to obtain the graph.

\num  Complete this statement: The function $f(x, y; p) = x^2 + p xy + 4 y^2$, which depends on the parameter $p$, has a minimum at the origin when \rule[-1pt]{30pt}{.3pt} and a minimax when \rule[-1pt]{30pt}{.3pt}\,.

\numa  Obtain the graph and a terraced density plot of the function $z = 3x^2 + 17 xy + 12 y^2$.  What is the shape of the graph?

\sub  What is the shape of the contours?  Indicate how the contours fit on the graph.

\numa  Obtain the graph and a terraced density plot of the function $z = 3x^2 + 7 xy + 12 y^2$.  What is the shape of the graph?

\sub  What is the shape of the contours?  Indicate how the contours fit on the graph.

\numa  Obtain the graph and a terraced density plot of the function $z = 3x^2 + 12 xy + 12 y^2$.  What is the shape of the graph?

\sub  What is the shape of the contours?  Indicate how the contours fit on the graph.


\num  Obtain the graph of $z = f(x, y) = xy$ on the domain $-3 \leq x \leq 3$, $-3 \leq y \leq 3$.  Does this function have a maximum or a minimum or a saddle point?  Where?

\numa  Continuation.  Sketch slices of the graph of $z = xy$ in the $y$-direction, for each of the values $x = -2$, $-1$, 0, 1, and 2.  What is the general shape of each of these slices?  

\sub  Repeat part (a), but make the five slices in the $x$-direction---that is, fix $y$ instead of $x$.  

\num  Continuation.  Show how the slices you obtained in the previous exercise fit (or appear) on the graph you obtained in the exercise just before that one.

\numa  Continuation.  Let $x = u + v$ and $y = u - v$.  Express $z$ in terms of $u$ and $v$, using the fact that $z = xy$.  Then obtain the graph of $z$ as a function of the new variables $u$ and $v$.  

\sub  What is the shape of the graph you just obtained?  Compare it to the graph of $z = xy$ you obtained earlier.  

\num  Can you draw a network of straight lines on the saddle surface $z = x^2 - y^2$?

\num  Obtain a terraced density plot of $z = xy$.  How do the contours of this plot fit on the graph of $z = xy$ you obtained in a previous exercise?

\num  The graphs of $z = x^2 + 5 xy + 10 y^2$ and $z = 3$ intersect in a curve.  What is the shape of that curve?

\num  The graphs of $z = x^2 + y^3$ and $z = 0$ intersect in a curve.  What is the shape of that curve?

\num  The graphs of $z = 2x - y$ and $z = .3 x + .8 y + 2.3$ intersect in a curve.  What is the shape of that curve?



\subsubsection{First integrals}
\vspace{-10pt}
\index{first integral}

\num  A hard spring\index{hard spring}\index{spring, hard} described by the dynamical system
	$$
	{dx \over dt} = v \qquad \quad {dv \over dt} = - cx - \beta x^3
	$$
has a first integral of the form
	$$
	E(x, v) = {\textstyle {1 \over 2}}c x^2 + 
		{\textstyle {1 \over 2}} \beta x^4 + 
		{\textstyle {1 \over 2}}v^2.
	$$
This is the {\bf energy}\index{energy} of the spring.  (See chapter~7, \S 3, especially exercise~13, page~447.)

\sub  Let $c = 16$ and $\beta = 1$.  Obtain the graph of $E(x, v)$ on a domain that has the origin at its center.  Locate all the minimum, maximum, and saddle points in this domain.  

\sub  What is the state of the spring (that is, its position $x$ and its velocity $v$) when it has minimum energy?  

\num  A soft spring\index{soft spring}\index{spring, soft} described by the dynamical system
	$$
	{dx \over dt} = v \qquad \quad {dv \over dt} = - {25x \over 1 + x^2}
	$$
has an energy integral of the form
	$$
	E(x, v) = {\textstyle {25 \over 2}} \ln(1 + x^2) + 
		{\textstyle {1 \over 2}}v^2.
	$$
(See exercise~16, page~448.)  

\sub  Obtain the graph of $E(x, v)$.  Experiment with different possibilities for the domain until you get a good representation.

\sub  Obtain a terraced density plot of $E(x, v)$ over the same domain you chose in part (a).  Compare the two representations of $E$.

\sub  Does the spring have a state of minimum energy?  If so, where is it?

\sub  Does the spring have a state of {\em maximum\/} energy?  Explain your answer.

\numa  {\bf The Lotka-Volterra equations}.\index{Lotka-Volterra equations}  According to exercise~33 of chapter~7, \S 3 (page~451), the function
	$$
	E(x, y) = .1 \ln{y} + .04 \ln{x} - .005 \, y -.004 \, x
	$$
is a first integral of the dynamical system
	\begin{eqnarray*}
	x' &=& .1 \, x - .005 \, xy	\\
	y' &=& .004 \, xy - .04 \, y.
	\end{eqnarray*}
Obtain the graph of $E$ on the domain $1\leq x \leq 50$, $1 \leq y \leq 50$.  (Why not enlarge the domain to $0 \leq x \leq 50$, $0 \leq y \leq 50$?)  

\sub  Find all maximum, minimum, and saddle points on this graph.  What is the connection between the maximum of $E$ and the equilibrium point of the dynamical system?

\sub  Obtain a contour plot of $E$ on the same domain as in part (a).  Compare the contours of $E$ and the trajectories of the dynamical system.  (This reveals a {\em conservation of ``energy''\/} for the solutions of the Lotka-Volterra equations.)

\numa  Continuation.  Here is another first integral of the same dynamical system as in the previous exercise:
	$$
	H(x, y) = x^{.04} y^{.1} / \exp(.005 \, y + .004 \, x).
	$$
Obtain the graph of $H$ and compare it to the graph of $E$ in the previous exercise.

\sub  Obtain a contour plot of $H$, and compare the contours to the trajectories of the dynamical system.


\end{document}
